\scanpage{134}
% archive foliation: 133
\[
\uline{f}^{m-n}:K_*(\uline{f}^{m})\qquad K_*(\uline{f}^{n})
\]
ou $m$ et $n$ sont des entiers rationnels, $m\quad n\quad0$, d'ou des \deleted{\unsure{en}}homomorphismes, notes par le meme symbole:
\[
K^*(\uline{f}^{n},M)\qquad K^*(\uline{f}^{m},M)
\]
et
\[
H^*(\uline{f}^{n},M)\qquad H^*(\uline{f}^{m},M)\qquad(m\quad n\quad0).
\]
Comme ces homomorphismes sont transitifs, on peut passer a la limite inductive dans ces homomorphismes, et poser
\[
C^*((\uline{f}^{\unsure{n}}),M)=\lim_n K^*(\uline{f}^{n},M)
\]
\[
\uline{H}^*((\uline{f}^{\unsure{n}}),M)\deleted{\unsure{*}}=H^*(C^*((\uline{f}),M)=\lim_n H^*(\uline{f}^{n},M).
\]
(N.B. dans ces notations, $(\uline{f})$ designe en fait l'ideal engendre dans $A$ par les $f_i$; on verra en effet ulterieurement que les objets introduits, \deleted{ne dependent} du moins $H^*((\uline{f}),M)$, ne depend que de cet ideal, et meme seulement de la topologie $\uline{m}$-adique qu'il definit, et non des $f_i$ eux-memes\unsure{.)}

Bien entendu, $C^*((\uline{f}),M)$ est un foncteur \deleted{exact} ($A$-lineaire exact) en $M$ et $H^*((\uline{f}),M)$ est un foncteur cohomologique en $M$.

Soit $\uline{f}$ comme ci-dessus, \deleted{\unsure{soit}} considere comme element du dual de $A^n$, soit $\uline{g}=(g_i)_{1\quad i\quad n}$ un element de $A^n$ lui-meme, on sait qu'on a alors, en designant par $e_{\uline{g}}$ l'operation de multiplication exterieure par $\uline{g}$ dans $A^n$, la \uline{formule d'homotopie}
\[
i_{\uline{f}}e_{\uline{g}}+e_{\uline{g}}i_{\uline{f}}=\langle\uline{g},\uline{f}\rangle\mathrm{id}.
\]
dans $\unsure{\Lambda}^{n}$, qui peut aussi s'ecrire, dans le complexe $K_*(\uline{f})$, comme
\[
de_{\uline{g}}+e_{\uline{g}}d=\uline{g},\uline{f}\quad\mathrm{id}.
\]
On en conclut la

\uline{Proposition 1}. Supposons que l'ideal $(\uline{f})$ engendre par les $f_i$ soit $A$. \unsure{Alors} le complexe $K_*(\uline{f})$ est homotopiquement equivalent a zero, et il en est \unsure{par} suite dememe des complexes $K_*(\uline{f},M)$ et $K^*(\uline{f},M)$ quels que soit $M$.
\margin{\unsure{Autre corollaire}: les $f_i$ annulent les $H_{\unsure{*}}(\uline{f},M)$, $H^{\unsure{*}}(\uline{f},M)$.}

\uline{Corollaire} S'il en est ainsi, on a $H^*(\uline{f},M)=0$ et $H^*((\uline{f}),M)=0$ pour \unsure{tout} $A$-module $M$.

Dans la description de $K^*(\uline{f},M)$, il faut aussi ajouter l'\unsure{interpretation} des elements de degre $p$ de ce complexe comme des fonctions alternees de $p$ indices variables entre $1$ et $n$; la formule du cobord est alors formellement celle d'un cobord en cohomologie de Cech.

2. \uline{Cohomologie de Cech d'un recouvrement ouvert}.

Soient $X$ un preschema, $A=\unsure{I|\Gamma}(X,\uline{O}_X)$, $\uline{f}=(f_i)_{1\quad i\quad n}$ un systeme de $n$ elements de $A$, $U_i=X_{f_i}$, $U=\quad U_i$, $\uline{U}$ le recouvrement de $U$ par les $U_i$. Soit $F$ un faisceau quasi-coherent sur $X$, posons
\[
M=\unsure{I|\Gamma}(X,F)
\]
(N.B. les notations precedentes devraient etre deployees de meme !)

